% Fragmento LaTeX generado desde docs/report/llm_model_section_final_es.md.
% Requiere, en el documento principal: \usepackage{booktabs}
% Recomendado: \usepackage{array}
% Para figuras: \usepackage{graphicx}
% Las figuras se esperan en una carpeta llamada figures/ junto al main.tex.
% Las figuras se insertan sin entorno flotante para respetar el orden exacto en Overleaf.

\makeatletter
\newcommand{\llmfixedfigure}[4]{%
\begin{center}
\refstepcounter{figure}\label{#4}%
\includegraphics[width=#1]{#2}
\par\vspace{2pt}
{\small\textbf{Figura~\thefigure.} #3}
\end{center}
}
\makeatother

\section{Sección para el report final: clasificación LLM de Planetary Boundaries}

\begin{quote}
Esta sección está pensada para integrarse en el informe final de UPV-EARTH dentro de los apartados de metodología, modelos, evaluación e iteración. Está escrita con foco en la parte de modelos de la rúbrica: elección técnica justificada, comparación de alternativas, protocolo de evaluación, mejora incremental, limitaciones y valor real del sistema.
\end{quote}

\subsection{1. Objetivo de modelado}

La parte LLM del proyecto UPV-EARTH aborda la clasificación automática de publicaciones científicas de la UPV respecto al marco de las 9 Planetary Boundaries (PBs). La unidad de análisis es el abstract limpio de cada publicación, enriquecido con metadatos auxiliares cuando están disponibles. La salida del sistema no es una etiqueta temática genérica, sino una decisión científica estructurada: una PB primaria, posibles PBs secundarias, una lista de PBs rechazadas y una explicación textual del razonamiento.

Desde el inicio se descartó tratar el problema como una simple búsqueda de palabras clave. En un corpus politécnico aparecen términos como \emph{water}, \emph{climate}, \emph{aerosol}, \emph{soil}, \emph{energy} o \emph{biodiversity} en contextos muy distintos. Un paper puede mencionar el cambio climático como motivación de fondo sin estudiar ninguna variable climática; otro puede ser de ingeniería hidráulica y usar \emph{water} sin tratar consumo global de agua dulce; otro puede hablar de sostenibilidad o política ambiental sin medir procesos biofísicos. El riesgo principal era por tanto el positivity bias o green hallucination: asignar PBs a cualquier documento con vocabulario ambiental superficial.

Por esa razón, el problema se formuló como una tarea de clasificación científica con abstención. El modelo debe poder responder \texttt{None} cuando no hay una PB activamente estudiada. En términos operativos, una PB se considera activa si el abstract mide, modela o impone experimentalmente una variable relacionada con esa frontera. Esta definición fue clave para separar tres niveles que suelen confundirse:

\begin{itemize}
\item \textbf{Mención contextual}: el abstract usa vocabulario ambiental como motivación, pero no lo estudia.
\item \textbf{Afinidad temática}: el abstract pertenece a un área próxima, como ingeniería, agricultura o planificación urbana, pero no cuantifica el proceso PB.
\item \textbf{Activación PB real}: el abstract mide, modela o manipula una variable que encaja con la definición y la lógica de activación de una PB.
\end{itemize}

Esta distinción convirtió el prompt en una pieza metodológica central, no en un simple envoltorio de inferencia.

\subsection{2. Datos, ground truth y protocolo de validación}

El desarrollo se realizó sobre el corpus curado definido en la entrega M2. La muestra inicial de 1,000 documentos se filtró mediante deduplicación y umbral mínimo de 500 caracteres en el abstract, produciendo una vista minable de 696 documentos con abstract suficientemente informativo. Para la fase final también existe un corpus institucional ampliado de 31,634 documentos limpios, pero la evaluación controlada de modelos se hizo sobre el subconjunto manualmente validado, porque es el único que permite medir rendimiento de forma honesta.

El ground truth está en \texttt{nlp/llm/outputs/ground\_truth/validacion\_real.csv}. Contiene 208 filas revisadas manualmente, 206 \texttt{doc\_ids} únicos y columnas \texttt{1stpb}, \texttt{2ndpb} y \texttt{3rdpb}. Al cruzarlo con los CSV actuales de inferencia se obtienen normalmente 150 filas evaluables y 149 \texttt{doc\_ids} únicos. La diferencia se debe a que el ground truth conserva algún documento duplicado con criterios de etiquetado distintos; esto se mantiene para ser consistente con los notebooks de análisis, pero se documenta como fuente de ruido.

Para evitar mezclar resultados antiguos con resultados finales, se usaron tres protocolos separados:

\begin{center}
\scriptsize
\setlength{\tabcolsep}{4pt}
\renewcommand{\arraystretch}{1.15}
\begin{tabular}{p{0.23\linewidth}p{0.16\linewidth}p{0.30\linewidth}p{0.23\linewidth}}
\toprule
Fuente de números & n usado & Para qué se usa & Nota de comparabilidad \\
\midrule
M2, Tabla 5.2 & Benchmark M2 & Comparar Llama 8B, Qwen 14B y Gemma 26B. & Usa \emph{Flexible Agreement} y \emph{Rigorousness}; no se mezcla con Top-1/Jaccard finales. \\
\texttt{problema\_few\_shot.pdf} & 150 documentos & Diagnosticar el primer Few-Shot contrastivo v1 frente a zero-shot. & Es la fuente correcta para las métricas históricas de v1 y para el fallo PB7 23\% $\rightarrow$ 0\%. \\
\texttt{comparacion\_fewshot\_v1\_v2\_v3\_v4\_principle.ipynb} & 147 filas, 146 \texttt{doc\_ids} & Comparar v1-v4 en igualdad de condiciones. & Excluye los 3 calibration cases de v4: \texttt{b9e1bf330baf}, \texttt{9153f4dcf3d6}, \texttt{f21bc832abc6}. \\
CSVs actuales de cascada & 150 filas, 149 \texttt{doc\_ids} & Evaluar la arquitectura multi-agente final. & Se recalcula con la misma lógica de métricas para que sea comparable con el notebook. \\
\bottomrule
\end{tabular}
\end{center}

El ground truth tiene naturaleza multi-etiqueta: un paper puede pertenecer a PB1 y PB4, o a PB9 y PB1, etc. Sin embargo, la evaluación principal se hizo sobre la primera PB humana (\texttt{1stpb}) frente a la PB primaria del modelo. Para no perder información multi-etiqueta, también se midieron métricas relajadas:

\begin{itemize}
\item \textbf{Top-1 estricto}: la PB primaria del modelo coincide exactamente con \texttt{1stpb}.
\item \textbf{Top-2 global}: la primera PB humana aparece como primaria o secundaria, incluyendo los aciertos \texttt{None -> None}.
\item \textbf{Hit@K PB-only}: la primera PB humana aparece como primaria o secundaria, pero calculado solo sobre documentos donde el humano sí asignó una PB.
\item \textbf{Exact Match de conjunto}: el conjunto de PBs predicho coincide con el conjunto humano.
\item \textbf{Jaccard}: solapamiento entre conjunto humano y conjunto predicho.
\item \textbf{Hamming Loss}: fracción de errores sobre las 9 PBs posibles.
\item \textbf{True Negative / Rigorousness}: proporción de documentos humanos \texttt{None} que el modelo rechaza correctamente.
\item \textbf{Positivity Bias}: proporción de documentos humanos \texttt{None} a los que el modelo asigna alguna PB.
\item \textbf{PB-only accuracy}: rendimiento sobre documentos donde sí hay PB humana.
\item \textbf{None-only accuracy}: rendimiento sobre documentos sin PB humana.
\end{itemize}

Esta batería de métricas evita optimizar una única cifra. En este problema, subir recall asignando PBs a casi todo no sirve, porque contamina el mapa institucional. Tampoco sirve rechazar casi todo como \texttt{None}, porque se pierden contribuciones reales. La evaluación debe capturar el equilibrio entre cobertura y rigor.

\begin{center}
\scriptsize
\setlength{\tabcolsep}{4pt}
\renewcommand{\arraystretch}{1.18}
\begin{tabular}{p{0.18\linewidth}p{0.28\linewidth}p{0.28\linewidth}p{0.18\linewidth}}
\toprule
Métrica & Definición operativa & Qué diagnostica & Riesgo si se optimiza sola \\
\midrule
Top-1 estricto & \texttt{pred\_primary} coincide con \texttt{1stpb}. & Calidad de la decisión principal que aparecería en un dashboard institucional. & Penaliza casos multi-PB donde el modelo acierta una PB secundaria razonable. \\
Top-2 global & \texttt{1stpb} aparece como primaria o secundaria, contando también \texttt{None -> None}. & Cobertura comparable en la tabla global. & Puede parecer alto si hay muchos \texttt{None} correctamente rechazados. \\
Hit@K PB-only & \texttt{1stpb} aparece como primaria o secundaria solo en documentos con PB humana. & Capacidad de no perder papers relevantes para revisión. & No mide falsos positivos sobre documentos irrelevantes. \\
Exact Match & El conjunto predicho coincide con el conjunto humano. & Precisión multi-etiqueta completa. & Muy severa con ground truth incompleto o discutible. \\
Jaccard & Intersección dividida por unión entre conjunto humano y predicho. & Solapamiento parcial en papers que tocan varias fronteras. & Puede parecer aceptable aunque la primaria esté mal. \\
None-only accuracy & Documentos humanos \texttt{None} rechazados como \texttt{None}. & Rigor y control de \emph{green hallucination}. & Favorece modelos excesivamente conservadores. \\
PB-only accuracy & Documentos con PB humana clasificados con la PB primaria correcta. & Recall real sobre contribuciones relevantes. & Favorece asignar PB a demasiados documentos. \\
Positivity bias & Fracción de \texttt{None} humanos convertidos en alguna PB. & Contaminación potencial del mapa institucional. & Si se minimiza demasiado, se pierden contribuciones reales. \\
\bottomrule
\end{tabular}
\end{center}

Esta tabla fue útil durante el desarrollo porque obligó a separar dos objetivos que a veces compiten: descubrir investigación relevante y no inflar artificialmente el impacto ambiental de la UPV. Por eso las versiones se compararon siempre con al menos una métrica de acierto, una métrica de rechazo y una métrica multi-etiqueta.

\subsection{3. Marco de referencia PB como base del prompt}

La pieza semántica central del sistema es \texttt{corpus\_PB/data/pb\_reference.csv}. Este archivo operacionaliza las 9 Planetary Boundaries en columnas utilizables por modelos y módulos deterministas:

\begin{itemize}
\item \texttt{pb\_code} y \texttt{pb\_name}: identificador y nombre de la frontera.
\item \texttt{short\_definition}: definición científica breve.
\item \texttt{core\_keywords}, \texttt{extended\_keywords}, \texttt{applied\_keywords\_upv}: vocabulario estratificado.
\item \texttt{activation\_logic}: condiciones bajo las cuales la PB debe activarse.
\item \texttt{exclusion\_notes}: condiciones que obligan a rechazar la PB aunque haya vocabulario superficial.
\item \texttt{source\_basis}: trazabilidad hacia documentos de referencia del proyecto.
\end{itemize}

La evolución metodológica importante fue pasar de \emph{definiciones de PB} a \emph{reglas de activación y exclusión}. Por ejemplo:

\begin{itemize}
\item PB1 no se activa por cualquier paper de energía; requiere emisiones, calentamiento, forzamiento radiativo, escenarios climáticos o impacto climático explícito.
\item PB2 no se activa por cualquier estudio marino; requiere acidificación, pH, carbonato, aragonito o efectos de acidez.
\item PB4 no se activa por agricultura en general; requiere nitrógeno, fósforo, eutrofización, fertilización o alteración de ciclos biogeoquímicos.
\item PB5 no se activa por calidad de agua o tratamiento de aguas; requiere cantidad, extracción, escasez, caudal, riego o presión sobre agua dulce.
\item PB7 no se activa por \emph{naturaleza} o \emph{conservación} genérica; requiere biodiversidad, especies, integridad ecosistémica, diversidad funcional/genética, pérdida de hábitat, etc.
\item PB8 exige sustancia o entidad novedosa concreta: plásticos, PFAS, pesticidas, fármacos, nanomateriales, sustancias persistentes o tóxicas.
\item PB9 exige partículas o aerosoles: PM2.5, PM10, AOD, black carbon, sulfate aerosols, aerosol-cloud interactions, etc.
\end{itemize}

Estas reglas se inyectan dinámicamente en cada prompt, de modo que la decisión del LLM queda anclada en el marco del proyecto y no solo en el conocimiento pre-entrenado del modelo.

\begin{center}
\scriptsize
\setlength{\tabcolsep}{4pt}
\renewcommand{\arraystretch}{1.18}
\begin{tabular}{p{0.08\linewidth}p{0.27\linewidth}p{0.28\linewidth}p{0.29\linewidth}}
\toprule
PB & Señales que activan & Señales que no bastan & Criterio práctico usado en prompts \\
\midrule
PB1 & Emisiones de GEI, calentamiento, forzamiento radiativo, escenarios climáticos, impactos climáticos modelados. & Energía, eficiencia, movilidad o sostenibilidad sin variable climática explícita. & Activar solo si el clima es variable medida, modelo, escenario o tratamiento, no si aparece como motivación general. \\
PB2 & pH oceánico, carbonatos, aragonito, acidificación marina, respuesta biogeoquímica a acidez. & Estudios marinos generales, biodiversidad marina sin química de carbonato. & Si la variable operacional es acidez/carbonato, PB2 puede superar a PB7 o PB1. \\
PB3 & Ozono estratosférico, ODS, CFCs, halones, agotamiento de ozono, radiación UV asociada. & Ozono troposférico o calidad del aire urbana sin capa estratosférica. & Exigir capa de ozono estratosférica o sustancias agotadoras; no confundir con smog fotoquímico. \\
PB4 & Nitrógeno, fósforo, fertilización, eutrofización, lixiviación, deposición de nutrientes. & Agricultura o productividad vegetal sin ciclo N/P. & Activar por flujo biogeoquímico medido, aunque el vector físico sea polvo, lluvia o escorrentía. \\
PB5 & Extracción, escasez, caudal, disponibilidad, riego, sequía hidrológica, presión sobre agua dulce. & Calidad de agua, tratamiento de aguas o contaminantes acuáticos sin cantidad/disponibilidad. & Separar agua como recurso físico de agua como medio contaminado. \\
PB6 & Cobertura del suelo, deforestación, urbanización, conversión agrícola, fragmentación, cambio de uso del suelo. & Localización rural/urbana mencionada solo como contexto. & Activar cuando el paper cuantifica transformación espacial o funcional del territorio. \\
PB7 & Biodiversidad, especies, abundancia, riqueza, hábitat, diversidad funcional/genética, integridad ecosistémica. & Naturaleza, paisaje, conservación o servicios ecosistémicos de forma genérica. & Activar si el paper mide respuesta biológica/ecológica, aunque el driver sea PB1 o uso del suelo. \\
PB8 & Plásticos, pesticidas, fármacos, PFAS, nanomateriales, metales pesados, compuestos persistentes/tóxicos. & Residuos, materiales o economía circular sin entidad contaminante concreta. & Exigir sustancia o clase de contaminante; ampliar vocabulario desde errores. \\
PB9 & PM2.5, PM10, AOD, black carbon, sulfatos, aerosoles, aerosol-cloud interactions. & Contaminación atmosférica genérica sin partículas/aerosoles. & Si se cuantifican partículas o aerosoles, PB9 no debe degradarse automáticamente a PB1. \\
\bottomrule
\end{tabular}
\end{center}

\subsection{4. Infraestructura de inferencia}

Los modelos se ejecutaron localmente mediante Ollama. Esto fue importante por tres motivos: coste cero por llamada, privacidad del corpus institucional y reproducibilidad en una VM controlada. Los runners usan Python, peticiones HTTP al endpoint local de Ollama y guardado incremental en CSV para no perder ejecuciones largas.

Las decisiones técnicas más relevantes fueron:

\begin{itemize}
\item \textbf{Modelos locales open-weight}: principalmente \texttt{qwen2.5:14b} para clasificación, \texttt{qwen2.5:3b} para extracción auxiliar y pruebas con \texttt{gemma4:26b} como modelo más grande.
\item \textbf{Temperatura 0.0}: se prioriza determinismo y comparabilidad entre versiones.
\item \textbf{JSON mode}: las salidas se fuerzan a JSON para poder parsear \texttt{primary\_pb}, \texttt{secondary\_pbs}, \texttt{rejected\_pbs}, \texttt{confidence} y \texttt{reasoning\_process}.
\item \textbf{Preflight de modelos}: el pipeline comprueba que Ollama responde y que los modelos están descargados antes de empezar, evitando runs corruptos por errores 404.
\item \textbf{Guardado incremental}: cada fila se escribe inmediatamente en CSV para tolerar caídas.
\item \textbf{Pipeline resumible}: los doc\_ids ya procesados se omiten salvo que se use \texttt{--no-resume}.
\item \textbf{Control de VRAM}: \texttt{keep\_alive=10m} evita que modelos cargados indefinidamente bloqueen la GPU.
\end{itemize}

La explicación generada por el LLM se guarda como parte del output. Esto no es solo \emph{texto bonito}: permite auditoría cualitativa de errores, identificación de patrones de fallo y revisión humana posterior.

\subsection{5. Evolución de modelos y prompts}

La parte más importante para la rúbrica es que no se probó un único modelo, sino una secuencia de hipótesis, errores medidos e iteraciones. La evolución puede resumirse como un proceso de calibración entre dos fallos opuestos:

\begin{itemize}
\item \textbf{Sobre-asignación}: el modelo ve PBs en cualquier paper con lenguaje ambiental.
\item \textbf{Infra-asignación}: el modelo se vuelve tan estricto que rechaza papers realmente relevantes.
\end{itemize}

\subsubsection{5.0 Mapa de iteraciones}

La siguiente tabla resume la lógica experimental. Cada versión cambió una hipótesis concreta del prompt o de la arquitectura; por tanto, las métricas no son una colección de pruebas aisladas, sino una secuencia de aprendizaje.

\begin{center}
\scriptsize
\setlength{\tabcolsep}{3pt}
\renewcommand{\arraystretch}{1.18}
\begin{tabular}{p{0.14\linewidth}p{0.21\linewidth}p{0.25\linewidth}p{0.25\linewidth}p{0.10\linewidth}}
\toprule
Versión & Hipótesis & Cambio técnico & Qué se observó & Decisión \\
\midrule
Zero-shot estricto & Un buen marco PB con exclusiones basta para clasificar. & Reglas completas, rol de evaluator, salida JSON y temperatura 0.0. & Buen rechazo de \texttt{None}, pero muchos falsos negativos en PB reales. & Mantener rigor, mejorar recall. \\
Few-shot v1 & Los ejemplos corregirán casos frontera. & Casos contrastivos dentro del prompt. & El modelo copió patrones y colapsó PB7; los ejemplos actuaron como plantillas rígidas. & Descartar esta familia de ejemplos. \\
V2 CoT guiado & Un checklist deductivo fuerza razonamiento completo. & Extracción de métricas, mapeo, exclusión, jerarquización. & Sube PB-only, pero explota el positivity bias. & Conservar checklist parcial, añadir filtro de foco. \\
V3 focus filter & Separar contexto de objeto operacional reducirá falsos positivos. & Verificación background-vs-focus y reglas para textos sociales/gobernanza. & Recupera rigor, pero pierde algunos PB reales. & Mantener principio, recalibrar casos difíciles. \\
V4 principle-driven & La primaria debe ser la variable medida/modelada, no la narrativa dominante. & Prompt XML, calibration cases no copiables, regla de objeto operacional. & Mejor single-model: 72.1\% Top-1 y 0.688 Jaccard. & Baseline final single-model. \\
Cascada multi-agente & Señales independientes aumentan trazabilidad y recall. & Extractor 3B, scorer léxico, juez 14B, critic asimétrico. & Top-1 cercano a v4, PB-only ligeramente superior y mucha más auditabilidad. & Opción final para human-in-the-loop. \\
\bottomrule
\end{tabular}
\end{center}

\subsubsection{5.1 Baseline zero-shot estricto}

El primer baseline fuerte fue \texttt{qwen2.5:14b} con un prompt zero-shot estructurado. Incluía rol de evaluador científico, reglas PB, exclusiones, extracción conceptual y salida JSON. El objetivo era medir si un LLM mediano podía aplicar el marco PB sin ejemplos.

La fuente primaria de esta ejecución es el CSV \texttt{nlp/llm/outputs/inferences/qwen2.5\_14b\_zeroshot.csv}. El resumen derivado está en \texttt{nlp/llm/outputs/studies/reporte\_zero\_shot.md}, generado el 2026-05-05. Para describir el baseline zero-shot de forma aislada se usan los valores del reporte; cuando se compara con v1-v4 en la tabla global, las métricas se recalculan desde el CSV para seguir la definición del notebook comparativo. La diferencia más visible es Jaccard: el reporte zero-shot da 0.476, mientras que la recomputación notebook-compatible da 0.626 porque trata los aciertos \texttt{None -> None} como solapamiento de conjunto.

Resultados del reporte zero-shot sobre 150 filas:

\begin{itemize}
\item Top-1: \textbf{63.3\%}
\item Hit@K PB-only: \textbf{56.6\%}
\item None-only accuracy / rigorousness: \textbf{84.3\%}
\item Positivity bias: \textbf{15.7\%}
\item Exact Match: \textbf{54.0\%}
\item Jaccard del reporte: \textbf{0.476}
\item Hamming loss: \textbf{0.064}
\end{itemize}

El comportamiento fue conservador. Rechazaba bien documentos irrelevantes, pero perdía demasiados papers con PB real. Los errores más frecuentes fueron \texttt{PB1 -> None}, \texttt{PB7 -> None} y \texttt{PB9 -> None}. Este baseline demostró que las reglas de exclusión funcionaban, pero también que el modelo se \emph{acobardaba} ante casos complejos.

\subsubsection{5.2 Few-shot v1: el fallo del ejemplo rígido}

La siguiente hipótesis fue añadir ejemplos few-shot para enseñar casos frontera. Intuitivamente, esto debería ayudar. Empíricamente ocurrió lo contrario.

La fuente correcta para el primer v1 no es el notebook final, sino \texttt{problema\_few\_shot.pdf}. Ese PDF analiza el archivo histórico \texttt{eval\_qwen2.5\_14b\_FEWSHOT\_208.csv} y compara zero-shot contra Few-Shot contrastivo sobre 150 documentos. Es importante separarlo de la recomputación apples-to-apples del notebook, porque el notebook reevalúa v1 sobre 147 filas para poder compararlo limpiamente con v4.

Resultados del PDF \texttt{problema\_few\_shot.pdf}:

\begin{itemize}
\item Exact Match: \textbf{49.3\%} frente a 54.0\% en zero-shot.
\item Top-1: \textbf{57.3\%} frente a 63.3\% en zero-shot.
\item Hit@K: \textbf{50.5\%} frente a 56.6\% en zero-shot.
\item Jaccard: \textbf{0.407} frente a 0.476 en zero-shot.
\item None-only accuracy / True Negative: \textbf{82.4\%} frente a 84.3\% en zero-shot.
\item Positivity bias: \textbf{17.6\%} frente a 15.7\% en zero-shot.
\item PB7 recall: \textbf{0\%}
\end{itemize}

El análisis cualitativo del PDF reveló el \emph{efecto loro}: el modelo reutilizaba frases del prompt en vez de razonar sobre el abstract. El documento cita una frase que aparecía repetida en la columna \texttt{llm\_reasoning}: \emph{Applying the Anti-Overfiltering clause... I must not over-reject}. En particular, internalizó reglas espurias, como tratar PB7 como un efecto secundario de PB1 en vez de permitir que biodiversidad fuese la frontera primaria. El PDF cuantifica este colapso como una caída del recall PB7 del 23\% al 0\%. También identifica la \emph{trampa social}: en 19 casos clasificados como \texttt{None}, el modelo justificó el rechazo diciendo que el texto era político o económico, incluso cuando había métricas biofísicas válidas. Esta iteración es valiosa precisamente porque muestra una decisión negativa justificada: no todo few-shot mejora un LLM; ejemplos mal diseñados pueden producir sobreajuste narrativo.

\subsubsection{5.3 Few-shot v2 / zero-shot CoT: checklist deductivo}

Para corregir v1, se eliminaron los ejemplos rígidos y se sustituyeron por un algoritmo deductivo obligatorio:

\begin{enumerate}
\item Extraer métricas físicas, químicas, biológicas o ecológicas explícitas.
\item Mapear cada métrica candidata a PBs.
\item Aplicar reglas de exclusión.
\item Jerarquizar primaria y secundarias.
\item Emitir JSON.
\end{enumerate}

Resultados en el mismo protocolo apples-to-apples de 147 filas:

\begin{itemize}
\item Top-1: \textbf{63.9\%}
\item PB-only accuracy: \textbf{69.1\%}
\item None-only accuracy: \textbf{54.0\%}
\item Positivity bias: \textbf{46.0\%}
\item Jaccard: \textbf{0.594}
\end{itemize}

V2 solucionó parte del problema de recall: encontraba muchos más papers con PB real. Pero lo hizo a costa de disparar falsos positivos. El positivity bias subió a casi la mitad de los documentos humanos \texttt{None}, lo que hacía la versión inaceptable para producción. Esta iteración enseñó que un checklist demasiado imperativo puede empujar al modelo a \emph{pescar} cualquier métrica posible.

\subsubsection{5.4 Few-shot v3: filtro background-vs-focus}

V3 mantuvo el razonamiento paso a paso, pero reforzó la distinción entre lo que el paper realmente mide/modela y lo que aparece solo como contexto. También añadió una reverificación para textos sociales, legales, de gobernanza, educación o software.

Resultados en el mismo protocolo apples-to-apples de 147 filas:

\begin{itemize}
\item Top-1: \textbf{64.6\%}
\item PB-only accuracy: \textbf{57.7\%}
\item None-only accuracy: \textbf{78.0\%}
\item Positivity bias: \textbf{22.0\%}
\item Jaccard: \textbf{0.596}
\end{itemize}

V3 corrigió gran parte del exceso de positividad de v2, pero volvió a perder recall en algunas clases. Es una versión más equilibrada que v2, aunque todavía por debajo del potencial del modelo.

\subsubsection{5.5 V4 principle-driven: objeto operacional y calibration cases}

La versión v4 fue el salto principal. Se reformuló el prompt con etiquetas XML y con un principio explícito: la PB primaria es la variable que el paper cuantifica, modela o impone experimentalmente, no la PB más mencionada retóricamente en la introducción.

La diferencia con v1 es importante. V4 no usa ejemplos como plantillas de respuesta, sino como calibration cases. Además, el prompt prohíbe reutilizar frases de esos casos. Los tres documentos usados como calibración se excluyen de la evaluación para no medir memorización.

El prompt introduce confusiones frecuentes:

\begin{itemize}
\item Aerosoles, PM y AOD son PB9, no PB1.
\item PB1 requiere clima activo: medido, modelado o impuesto como escenario/tratamiento.
\item Papers de gestión de agua con clima solo como motivación son PB5, no PB1.
\item Respuestas ecológicas y biodiversidad pueden ser PB7 primaria.
\item Nutrientes N/P apuntan a PB4.
\end{itemize}

Resultados sobre 147 filas cruzadas:

\begin{itemize}
\item Top-1: \textbf{72.1\%}
\item PB-only accuracy: \textbf{70.1\%}
\item None-only accuracy: \textbf{76.0\%}
\item Top-2 global: \textbf{74.8\%}
\item Hit@K PB-only: \textbf{74.2\%}
\item Exact Match de conjunto: \textbf{60.5\%}
\item Jaccard: \textbf{0.688}
\item Hamming loss: \textbf{0.061}
\item Positivity bias: \textbf{24.0\%}
\item Macro-F1: \textbf{0.679}
\item Weighted-F1: \textbf{0.718}
\item Tiempo medio por inferencia: aproximadamente \textbf{5 s/documento} en el análisis comparativo.
\end{itemize}

V4 fue la mejor versión single-model. Su mejora no viene de sacrificar una métrica por otra: gana claramente en las métricas de acierto y se mantiene razonablemente alto en rechazo de \texttt{None}. Frente a v3, el análisis pareado mostró 15 documentos donde v4 acierta y v3 falla, frente a solo 4 donde ocurre lo contrario. Por tanto, la mejora no es un artefacto agregado, sino una ganancia real a nivel documento.

\subsubsection{5.6 Anatomía de prompts}

Para que la evolución no quede como una explicación abstracta, se incluyen aquí extractos literales de los prompts reales. La fuente no son reconstrucciones, sino los runners/notebooks del repositorio: \texttt{qwen\_fewshot\_v2.py}, \texttt{qwen\_fewshot\_v3.ipynb}, \texttt{qwen\_fewshot\_v4.ipynb} y \texttt{pipeline\_agentes.py}. En los bloques se omiten solo las partes dinámicas demasiado largas, como \texttt{\{pb\_rules\}} y \texttt{\{abstract\_text\}}, porque se rellenan en tiempo de ejecución desde \texttt{pb\_reference.csv} y el abstract de cada paper.

\paragraph{Prompt v2 real: checklist deductivo}

{\footnotesize
\begin{verbatim}
You are a strict scientific evaluator analyzing research abstracts
from the Universitat Politecnica de Valencia (UPV) against the
Planetary Boundaries (PBs) framework.

Your task: identify the Planetary Boundaries that the research
actually measures or models.

### YOUR ANALYTICAL ALGORITHM (mandatory deductive checklist)

Execute these steps internally, in order, ON THE REAL TEXT BELOW.
Do not invent metrics that are not in the abstract. Do not reuse
rhetoric from these instructions in your output.

Step 1 - BIOPHYSICAL EXTRACTION
   List every explicit physical, chemical, biological, ecological or
   environmental metric, variable, model, sample, dataset or measurement
   that appears in the abstract. Quote them.
   - If the list is EMPTY (... no biophysical observation), the verdict
     is "None" and you skip to Step 5.

Step 2 - CANDIDATE MAPPING
   For EACH extracted metric, map it to the PB whose Core Definition /
   Activation Logic it triggers.
   - Biosphere/ecological responses (...) -> PB7 is a legitimate PRIMARY.
   - Particulate matter, PM2.5, PM10, aerosols, AOD, black carbon -> PB9,
     NOT PB1.
   - "Climate" or "warming" alone is NOT enough for PB1...

Step 3 - EXCLUSION CHECK
   For every candidate PB from Step 2, apply its EXCLUSION RULE.

Step 4 - HIERARCHY
   Select as `primary_pb` the boundary that is the dominant analytical focus.

### OUTPUT FORMAT
Output ONLY valid JSON matching this exact schema, no markdown, no extra text:
{
  "reasoning_process": "Step1 metrics found: [...]. Step2 candidates: [...].",
  "primary_pb": {"pb_code": "PBX", "confidence": "High/Medium/Low"} or null,
  "secondary_pbs": [{"pb_code": "PBY", "confidence": "High/Medium/Low"}],
  "rejected_pbs": ["PBZ"]
}
\end{verbatim}
}

La intuición de v2 era correcta: obligar al modelo a nombrar variables antes de elegir una PB. El problema fue que el imperativo de extracción producía una presión excesiva a encontrar algo clasificable. Si el abstract contenía cualquier métrica ambiental débil, el modelo tendía a convertirla en PB activa.

\paragraph{Prompt v3 real: filtro background-vs-focus}

{\footnotesize
\begin{verbatim}
Step 1 - BIOPHYSICAL EXTRACTION & VALIDATION
   List the physical, chemical, biological, or ecological metrics that this
   specific research ACTUALLY MEASURES, MODELS, or ANALYZES. Quote them.
   - CRITICAL DISTINCTION: Differentiate between "background context" and
     "research focus". If the abstract merely mentions "climate change",
     "warming", or "emissions" as the background reason for a purely
     economic, social, legal, or governance study, IT DOES NOT COUNT.
   - OPERATIONAL CRITERION: A metric counts as "actually analyzed" only if
     the abstract reports a number, a comparison, a model output, or a
     measurement methodology around it.
   - If the list of *actually analyzed* biophysical metrics is EMPTY, the
     verdict is "None" and you skip to Step 5.

Step 3 - EXCLUSION CHECK
   For every candidate PB from Step 2, apply its EXCLUSION RULE.
   - Re-verify: Is the paper purely social science, policy, or theoretical
     computing? If yes, and the biophysical terms were just background
     context, discard the PB to prevent positivity bias.
\end{verbatim}
}

V3 corrigió el exceso de v2 introduciendo una pregunta previa: qué hace realmente el paper. Esta regla fue especialmente importante en abstracts de gobernanza, educación, diseño urbano o sostenibilidad, donde el lenguaje ambiental puede ser intenso aunque no haya una frontera planetaria medida.

\paragraph{Prompt v4 real: principle-driven + XML}

{\footnotesize
\begin{verbatim}
<system_role>
You are an expert scientific evaluator analyzing research abstracts from the
Universitat Politecnica de Valencia (UPV) against the Planetary Boundaries
(PBs) framework. Your judgment matters more than rigid rule-following.
Use the provided framework as a reasoning aid, not a script.
</system_role>

<instructions>
1. Identify the *operational object*: Determine the exact variable being
   measured, modelled, or experimentally manipulated.
2. Match with PB: The *primary* PB is the one being analytically resolved,
   not the one mentioned in the introduction.
3. Apply EXCLUSION RULES...
4. Check common confusions:
   - Aerosols / PM / AOD belong to PB9, not PB1.
   - "Climate" or "warming" alone is not enough for PB1.
   - Biodiversity / ecosystem responses are a legitimate primary candidate
     for PB7.
   - Nutrients (N, P) point to PB4.
   - Water quantity/quality point to PB5.
</instructions>

<calibration_cases>
These cases share one principle: *the primary PB is whatever variable the
paper quantifies or models*, regardless of the introduction's rhetorical
framing. Use them ONLY to calibrate your judgment, NOT to mimic their wording.

- Case A: A paper studies the central carbon metabolism of anammox bacteria
  using 13C isotope tracing... Primary PB: PB4.
- Case B: A paper verifies the WRF forecast model... Climate is motivational
  backdrop, not a measured variable. Primary PB: None.
- Case C: A paper quantifies detection data for 24 mammal species in response
  to human footprint... Primary PB: PB7.
</calibration_cases>

<constraints>
- DO NOT reuse phrases from the calibration cases. Your abstract deserves
  a fresh reading.
- Write your own reasoning, citing the abstract's own terms.
- DO NOT output any markdown formatting (like ```json).
</constraints>
\end{verbatim}
}

La diferencia crítica de v4 es que los ejemplos ya no enseñan frases, sino límites de decisión. Esto redujo el \emph{efecto loro} de v1 y evitó que el modelo tratase los calibration cases como plantillas. También hizo explícita una regla científica difícil: una PB puede aparecer como causa global, pero la etiqueta primaria debe seguir el objeto operacional del paper.

\subsection{6. Comparativa global de sistemas LLM}

Tabla recomputada desde los CSV actuales de \texttt{nlp/llm/outputs/inferences/}, el notebook \texttt{comparacion\_fewshot\_v1\_v2\_v3\_v4\_principle.ipynb} y los outputs actuales de \texttt{nlp/llm/outputs/pipeline\_cascada/}. En v1-v4 se usa el protocolo del notebook, por eso aparecen 147 filas; en zero-shot y cascada se usan las 150 filas que cruzan con el ground truth actual.

\begin{center}
\scriptsize
\setlength{\tabcolsep}{3pt}
\renewcommand{\arraystretch}{1.15}
\begin{tabular}{p{0.22\linewidth}p{0.05\linewidth}p{0.065\linewidth}p{0.065\linewidth}p{0.065\linewidth}p{0.065\linewidth}p{0.065\linewidth}p{0.065\linewidth}p{0.065\linewidth}p{0.065\linewidth}}
\toprule
Sistema & n eval & Top-1 & PB-only & None-only & Hit@K PB-only & Exact set & Jaccard & Hamming $\downarrow$ & Positivity bias $\downarrow$ \\
\midrule
Qwen 14B zero-shot & 150 & 63.3\% & 52.5\% & 84.3\% & 56.6\% & 56.0\% & 0.626 & 0.064 & 15.7\% \\
Qwen 14B v1 contrastivo & 147 & 58.5\% & 45.4\% & \textbf{84.0\%} & 50.5\% & 53.1\% & 0.581 & 0.067 & 16.0\% \\
Qwen 14B v2 CoT guiado & 147 & 63.9\% & 69.1\% & 54.0\% & 72.2\% & 50.3\% & 0.594 & 0.071 & 46.0\% \\
Qwen 14B v3 focus filter & 147 & 64.6\% & 57.7\% & 78.0\% & 61.9\% & 52.4\% & 0.596 & 0.073 & 22.0\% \\
Qwen 14B v4 principle & 147 & \textbf{72.1\%} & 70.1\% & 76.0\% & 74.2\% & \textbf{60.5\%} & \textbf{0.688} & \textbf{0.061} & 24.0\% \\
Cascada multi-agente base & 150 & 64.0\% & 60.6\% & 70.6\% & 66.7\% & 54.0\% & 0.623 & 0.073 & 29.4\% \\
Cascada + critic inicial & 150 & 65.3\% & 62.6\% & 70.6\% & 68.7\% & 54.7\% & 0.633 & 0.072 & 29.4\% \\
Cascada v3 lean + critic (último run) & 150 & 71.3\% & \textbf{71.7\%} & 70.6\% & \textbf{74.7\%} & 58.0\% & 0.666 & 0.064 & 29.4\% \\
\bottomrule
\end{tabular}
\end{center}

La mejor variante estrictamente por Top-1 sigue siendo Qwen v4 single-model, con 72.1\%. La última cascada queda muy cerca, con 71.3\%, mejora ligeramente el PB-only accuracy (71.7\% frente a 70.1\%) y queda prácticamente empatada en Hit@K PB-only (74.7\% frente a 74.2\%). El log del runner de la cascada informa 120/150 = 80.0\% con una comprobación de solapamiento más laxa; para la tabla se usa la definición comparable del notebook, basada en documentos con PB humana. Esta diferencia explica la recomendación final:

\begin{itemize}
\item Para \textbf{batch labeling} donde importa maximizar exactitud primaria, v4 principle-driven es la opción más limpia.
\item Para \textbf{sistema auditable y human-in-the-loop}, la cascada es más defendible porque produce trazas intermedias: extracción del agente pequeño, evidencias léxicas, razonamiento del juez y posible revisión del critic.
\end{itemize}

\begin{center}
\scriptsize
\setlength{\tabcolsep}{5pt}
\renewcommand{\arraystretch}{1.2}
\begin{tabular}{p{0.22\linewidth}p{0.30\linewidth}p{0.30\linewidth}}
\toprule
Sistema clave & Top-1 estricto & Positivity bias \\
\midrule
Zero-shot & \rule{3.08cm}{0.8ex} 63.3\% & \rule{0.85cm}{0.8ex} 15.7\% \\
v2 CoT & \rule{3.11cm}{0.8ex} 63.9\% & \rule{2.50cm}{0.8ex} 46.0\% \\
v3 focus filter & \rule{3.14cm}{0.8ex} 64.6\% & \rule{1.20cm}{0.8ex} 22.0\% \\
v4 principle & \rule{3.50cm}{0.8ex} 72.1\% & \rule{1.30cm}{0.8ex} 24.0\% \\
Cascada v3 lean & \rule{3.46cm}{0.8ex} 71.3\% & \rule{1.60cm}{0.8ex} 29.4\% \\
\bottomrule
\end{tabular}
\end{center}

Este gráfico de barras en LaTeX deja ver visualmente la conclusión: v4 no gana por volverse más permisivo como v2, sino por subir Top-1 manteniendo el sesgo de positividad en una zona razonable.

\llmfixedfigure{0.92\linewidth}{figures/fig3_compare_systems.png}{Comparación visual de sistemas sobre el conjunto de validación. La figura resume la tensión entre accuracy global, rendimiento sobre documentos con PB real y capacidad de rechazo en documentos \texttt{None}.}{fig:llm-system-comparison}

\llmfixedfigure{0.90\linewidth}{figures/fig6_pipeline_vs_v4.png}{Comparación documento a documento entre el baseline v4 y la cascada. Esta visualización es útil para defender que el análisis no se limita a una media agregada: muestra qué sistema acierta o falla en cada paper.}{fig:pipeline-vs-v4}

\llmfixedfigure{0.95\linewidth}{figures/fig5_per_class_f1.png}{Métricas por clase de la cascada. La variación entre PBs muestra que el rendimiento global no debe interpretarse de forma homogénea: PB9 es más estable por su vocabulario físico distintivo, mientras que PB5/PB6/PB8 sufren por bajo soporte o vocabulario incompleto.}{fig:per-class-f1}

\subsection{7. Modelos grandes: Gemma 26B y selección de Qwen}

Durante M2 se compararon escalas de modelo: Llama 3.1 8B, Qwen 2.5 14B y Gemma 4 26B. Esta comparación pertenece al protocolo antiguo de M2, por lo que no debe mezclarse directamente con la tabla final v1-v4. Aun así, es relevante para justificar la elección de arquitectura porque mide coste, rigurosidad y comportamiento relativo entre modelos bajo el mismo prompt base.

\begin{center}
\scriptsize
\setlength{\tabcolsep}{5pt}
\renewcommand{\arraystretch}{1.15}
\begin{tabular}{p{0.23\linewidth}p{0.18\linewidth}p{0.18\linewidth}p{0.18\linewidth}p{0.15\linewidth}}
\toprule
Modelo M2 & Flexible Agreement & Rigorousness & Positivity Bias & Tiempo medio \\
\midrule
Llama 3.1 8B & 60.27\% & 0.00\% & 10.96\% & 3.08 s/doc \\
Qwen 2.5 14B & 52.05\% & 38.36\% & 0.00\% & 6.12 s/doc \\
Gemma 4 26B & 47.95\% & 39.73\% & 0.00\% & 22.90 s/doc \\
\bottomrule
\end{tabular}
\end{center}

\begin{center}
\scriptsize
\setlength{\tabcolsep}{5pt}
\renewcommand{\arraystretch}{1.2}
\begin{tabular}{p{0.22\linewidth}p{0.32\linewidth}p{0.32\linewidth}}
\toprule
Modelo M2 & Agreement visual & Latencia visual \\
\midrule
Llama 3.1 8B & \rule{3.50cm}{0.8ex} 60.27\% & \rule{0.51cm}{0.8ex} 3.08 s \\
Qwen 2.5 14B & \rule{3.02cm}{0.8ex} 52.05\% & \rule{1.02cm}{0.8ex} 6.12 s \\
Gemma 4 26B & \rule{2.78cm}{0.8ex} 47.95\% & \rule{3.80cm}{0.8ex} 22.90 s \\
\bottomrule
\end{tabular}
\end{center}

La lectura correcta de M2 no era que Llama fuese mejor modelo científico por tener más agreement flexible. De hecho, Llama tenía 0.00\% de rigorousness: no filtraba bien los documentos irrelevantes. Gemma era mejor como auditor conceptual: detectaba rutas químicas o biofísicas más profundas, por ejemplo relacionando SO2 con sulfatos particulados o distinguiendo un estudio literario sobre crisis climática de un paper climático real. Sin embargo, Gemma tenía dos problemas para producción:

\begin{itemize}
\item \textbf{Latencia}: en la evaluación M2 rondaba 22.9 s/documento frente a unos 6.1 s/documento de Qwen.
\item \textbf{Robustez operativa}: los runners finales muestran que Gemma requería más cuidado con formato chat y JSON schema para evitar outputs inválidos o tokens basura.
\end{itemize}

Además, el CSV final \texttt{gemma4\_26b\_fewshot\_v4\_principle.csv} contiene solo 2 filas de datos, por lo que no es válido como benchmark final. Por eso los números cuantitativos de Gemma/Llama se citan desde M2 y se etiquetan explícitamente como comparación histórica de arquitectura, mientras que las métricas finales se toman del notebook v1-v4 y de los CSV actuales de cascada.

Por esto, Gemma se mantuvo como referencia cualitativa y posible fuente de destilación futura, pero Qwen 2.5 14B fue seleccionado como compromiso operativo: suficientemente razonador, mucho más rápido y estable en Ollama.

Esta decisión no fue solo \emph{elegir el que más acierta}. Fue una decisión de sistema: el objetivo final es procesar miles o decenas de miles de publicaciones, no solo ganar una tabla pequeña. En ese contexto, una arquitectura mediana con prompts robustos y auditoría posterior es más viable que un modelo grande lento y difícil de estabilizar.

\subsection{8. Arquitectura multi-agente final}

La iteración final exploró una cascada multi-agente implementada en \texttt{nlp/llm/runners/pipeline\_agentes.py}. La idea no era reemplazar el prompt v4, sino enriquecerlo con señales independientes.

\llmfixedfigure{0.98\linewidth}{figures/fig1_architecture.png}{Arquitectura de la cascada multi-agente. La clave de diseño es que cada bloque tiene una responsabilidad acotada: extracción, evidencia léxica, juicio científico y crítica asimétrica.}{fig:cascade-architecture}

\llmfixedfigure{0.98\linewidth}{figures/fig2_data_flow.png}{Flujo de datos por documento. Las señales auxiliares se pasan al juez como evidencia no autoritativa para evitar que sustituyan la lectura directa del abstract.}{fig:cascade-flow}

\begin{center}
\scriptsize
\setlength{\tabcolsep}{4pt}
\renewcommand{\arraystretch}{1.18}
\begin{tabular}{p{0.14\linewidth}p{0.22\linewidth}p{0.24\linewidth}p{0.30\linewidth}}
\toprule
Bloque & Entrada & Salida & Control de fallo \\
\midrule
Agente 1 extractor & Título y abstract. & Listas de especies, métricas, observaciones, metodología y disciplina. & Grounding literal: se descartan strings que no aparecen en el texto. \\
Scorer léxico & Texto enriquecido y \texttt{pb\_reference.csv}. & Ranking de PBs candidatas con keywords activadoras. & Pesos simples y trazables; no decide por sí solo. \\
Router & Extracción y ranking. & Fast-skip a \texttt{None} o llamada al juez. & Solo salta cuando extractor y scorer coinciden en irrelevancia fuerte. \\
Agente 3 juez & Abstract, reglas PB, casos v4 y señales auxiliares. & PB primaria, secundarias, rechazadas, confianza y razonamiento. & Señales auxiliares marcadas como falibles para reducir anclaje. \\
Agente 4 critic & Casos \texttt{None} con candidatos léxicos. & Override limitado a \texttt{None -> PBx}. & Crítica asimétrica: no puede degradar aciertos PB ni inventar candidatos. \\
\bottomrule
\end{tabular}
\end{center}

\subsubsection{8.1 Agente 1: extractor estructurado (\texttt{qwen2.5:3b})}

El agente pequeño no decide la PB. Extrae campos estructurados:

\begin{itemize}
\item \texttt{chemical\_species}
\item \texttt{physical\_metrics}
\item \texttt{biological\_observations}
\item \texttt{methodology}: \texttt{measured}, \texttt{modelled}, \texttt{mentioned} o \texttt{none}
\item \texttt{disciplinary\_frame}: ingeniería, earth sciences, biología, química, social sciences, economics, law/policy, education u other
\end{itemize}

Para evitar alucinaciones, el prompt exige que cada string extraído aparezca literalmente en el título o abstract. Además, el código aplica un filtro posterior de grounding: cualquier item que no aparezca en el texto se descarta. Esto responde directamente a un fallo observado en prompts anteriores: los modelos pequeños copian ejemplos o inventan variables si se les deja demasiado margen.

\paragraph{Prompt del Agente 1}

{\footnotesize
\begin{verbatim}
<task>
Extract only explicit evidence from the title and abstract.
Do not decide the Planetary Boundary.
</task>

<hard_constraints>
- Every extracted string must be a direct substring of the input text.
- If the paper only mentions an environmental topic as motivation,
  put the methodology as mentioned.
- If no physical, chemical, biological or ecological evidence exists,
  return empty lists.
</hard_constraints>

<json_schema>
{
  "chemical_species": [],
  "physical_metrics": [],
  "biological_observations": [],
  "methodology": "measured|modelled|mentioned|none",
  "disciplinary_frame": "engineering|earth_sciences|biology|chemistry|social_sciences|other"
}
</json_schema>
\end{verbatim}
}

Este prompt es deliberadamente estrecho. No se le pide al modelo pequeño que razone sobre PBs porque en pruebas preliminares tendía a sobreinterpretar. Su valor es producir una capa de evidencia barata, no una decisión semántica final.

\subsubsection{8.2 Scorer determinista por vocabulario}

El segundo módulo no usa LLM. Calcula overlap entre \texttt{title + clean\_abstract + top\_terms\_no\_stopwords} y el vocabulario de \texttt{pb\_reference.csv}, con pesos:

\begin{itemize}
\item \texttt{core\_keywords}: peso 3
\item \texttt{extended\_keywords}: peso 2
\item \texttt{applied\_keywords\_upv}: peso 1
\end{itemize}

Cada keyword cuenta como máximo una vez y se descartan términos de menos de 4 caracteres. El resultado es una lista de candidatos tipo \texttt{PB1(6); PB5(3)} junto con las keywords que activaron la señal.

Este módulo es transparente y barato. Su limitación es la cobertura léxica: si PB8 no contiene heavy metals, un paper sobre Pb, Cu o Mn puede quedar sin candidato aunque sea relevante. Por eso el scorer se usa como pista, no como verdad.

\subsubsection{8.3 Router de consenso}

El router solo evita llamar al modelo grande cuando hay consenso fuerte de irrelevancia:

\begin{enumerate}
\item El Agente 1 devuelve cero items y un marco no biofísico.
\item El scorer no encuentra candidatos con score suficiente.
\end{enumerate}

En el último run el fast-skip se activó en 5 casos y no introdujo errores, pero su impacto en coste es pequeño. Su valor principal es demostrar que el sistema puede incorporar rutas baratas sin sacrificar calidad.

\subsubsection{8.4 Agente 3: juez experto (\texttt{qwen2.5:14b})}

El Agente 3 es el clasificador principal. Recibe:

\begin{itemize}
\item El abstract crudo.
\item Las reglas PB completas.
\item Los calibration cases v4.
\item La extracción del Agente 1.
\item El ranking del scorer.
\end{itemize}

En la versión final, las señales auxiliares se colocan después del abstract y se etiquetan explícitamente como no autoritativas. Esto fue una corrección importante: en runs previos, el 14B se anclaba demasiado en \texttt{kw\_top} o en la salida del extractor. Si \texttt{kw\_top} empataba PB1 y PB2, el modelo podía elegir PB1 por posición, aunque el abstract tratase ocean acidification. La versión final avisa al modelo de que los empates no son significativos y de que debe resolver leyendo el abstract.

\paragraph{Prompt del Agente 3}

{\footnotesize
\begin{verbatim}
<abstract_first>
Read the title and abstract before looking at auxiliary signals.
The abstract is authoritative. Auxiliary signals are fallible hints.
</abstract_first>

<pb_rules>
Use the PB reference definitions, activation logic and exclusions.
Apply the operational object principle from v4.
</pb_rules>

<auxiliary_evidence>
Agent 1 extracted evidence: may be incomplete or wrong.
Keyword scorer candidates: may reflect vocabulary overlap only.
Ties in keyword score are not meaningful rankings.
</auxiliary_evidence>

<decision_policy>
Choose None when no PB variable is actively measured, modelled
or imposed. Choose a PB when evidence is explicit even if the
keyword scorer is weak.
</decision_policy>
\end{verbatim}
}

\subsubsection{8.5 Agente 4: critic asimétrico}

El critic se invoca solo cuando:

\begin{itemize}
\item El Agente 3 responde \texttt{None}.
\item El scorer sí encontró al menos una PB candidata.
\end{itemize}

El critic solo puede cambiar \texttt{None -> PBx}. No puede degradar \texttt{PBx -> None} ni inventar una PB fuera de los candidatos. Esta asimetría se diseñó a partir de un diagnóstico empírico: una parte de los errores eran falsos negativos, pero no queríamos abrir una segunda pasada que destruyese aciertos o aumentase aún más el positivity bias.

\paragraph{Prompt del critic}

{\footnotesize
\begin{verbatim}
<critic_scope>
The main classifier returned None, but the keyword scorer found
candidate PBs. Review only whether this is a false negative.
</critic_scope>

<allowed_actions>
- keep None,
- override None to one candidate PB if explicit evidence exists.
</allowed_actions>

<forbidden_actions>
- do not change one PB into another PB,
- do not remove an existing PB prediction,
- do not choose a PB outside the candidate list,
- do not override from background motivation alone.
</forbidden_actions>

<output_json>
{
  "action": "keep_none|override",
  "override_pb": "PBx or None",
  "reason": "short evidence-based explanation"
}
</output_json>
\end{verbatim}
}

En el último run (\texttt{pipeline\_20260510\_192035.log}) la cascada procesó 149 doc\_ids únicos, generó 150 filas de evaluación al cruzar con el GT, activó fast-skip 5 veces, llamó al juez grande 144 veces, invocó el critic 5 veces y produjo 2 overrides. El resultado final fue:

\begin{itemize}
\item Accuracy vs GT 1st PB: \textbf{107/150 = 71.3\%}
\item Top-2 global recomputado con la definición del notebook: \textbf{110/150 = 73.3\%}
\item Hit@K PB-only: \textbf{74/99 = 74.7\%}
\item Solapamiento amplio informado por el log del runner: \textbf{120/150 = 80.0\%}
\end{itemize}

La cascada no supera claramente a v4 en Top-1 ni en Top-2 global, pero queda muy cerca y ofrece un diseño más auditable. También mejora levemente el rendimiento PB-only frente a v4 (71.7\% frente a 70.1\%). Este es un resultado honesto y defendible: la arquitectura multi-agente aporta trazabilidad y capacidad de revisión, aunque no siempre aumenta la métrica primaria.

\llmfixedfigure{0.82\linewidth}{figures/fig7_critic_impact.png}{Impacto del critic asimétrico. La figura separa los casos conservados como \texttt{None} y los overrides, mostrando que la crítica se usa como mecanismo puntual de recuperación de falsos negativos, no como segunda clasificación general.}{fig:critic-impact}

\subsection{9. Diagnóstico de errores}

El análisis de errores reveló patrones útiles para futuras iteraciones.

\llmfixedfigure{0.72\linewidth}{figures/fig4_confusion_matrix.png}{Matriz de confusión de la cascada final. Los errores fuera de la diagonal permiten distinguir tres familias: \texttt{PBx -> None}, \texttt{None -> PBx} y confusiones entre PBs activas.}{fig:confusion-matrix}

\begin{center}
\scriptsize
\setlength{\tabcolsep}{4pt}
\renewcommand{\arraystretch}{1.18}
\begin{tabular}{p{0.18\linewidth}p{0.17\linewidth}p{0.25\linewidth}p{0.32\linewidth}}
\toprule
Familia de error & Frecuencia aproximada & Causa principal & Mitigación propuesta \\
\midrule
\texttt{PBx -> None} & 15 casos & Prompt demasiado conservador o scorer sin vocabulario suficiente. & Critic más amplio y expansión de \texttt{pb\_reference.csv} desde falsos negativos. \\
\texttt{None -> PBx} & 15 casos & Positivity bias o ground truth humano discutible en papers con aerosoles/contaminantes reales. & Reanotación doble y etiqueta \texttt{Uncertain} para discrepancias científicas. \\
\texttt{PBx -> PBy} & 24 casos & Confusión entre driver global y variable medida; anclaje en señales auxiliares. & Reforzar objeto operacional y bajar autoridad de \texttt{kw\_top}. \\
Errores de parseo & No dominante & JSON mode y validación redujeron fallos de formato. & Mantener contratos JSON y tests de parser antes de runs largos. \\
\bottomrule
\end{tabular}
\end{center}

\llmfixedfigure{0.92\linewidth}{figures/fig8_failure_taxonomy.png}{Taxonomía de fallos observada en el análisis de la cascada. El valor metodológico de la taxonomía es que convierte errores agregados en acciones concretas de mejora.}{fig:failure-taxonomy}

Además de contar errores, se revisaron documentos concretos. Esta revisión cualitativa fue imprescindible porque algunas discrepancias no son simples fallos del modelo, sino ambigüedades del propio marco de anotación.

\begin{center}
\scriptsize
\setlength{\tabcolsep}{3pt}
\renewcommand{\arraystretch}{1.15}
\begin{tabular}{p{0.12\linewidth}p{0.20\linewidth}p{0.18\linewidth}p{0.20\linewidth}p{0.22\linewidth}}
\toprule
Doc & Fenómeno observado & Predicción problemática & Lectura humana/técnica & Lección de prompt \\
\midrule
\texttt{0866b31f} & Estudio literario con crisis climática como tema cultural. & PB1 en modelos menos estrictos. & \texttt{None}: no mide clima ni variable biofísica. & Excluir humanities/policy si no hay métrica ambiental. \\
\texttt{733ba5fc} & pH oceánico y química de acidificación. & PB1 por contexto climático. & PB2 como objeto operacional. & Separar driver climático de variable química medida. \\
\texttt{260cc670} & Deposición atmosférica asociada a nutrientes. & PB9 por vector aerosol. & PB4 si el resultado principal son flujos N/P. & Distinguir vector físico de ciclo biogeoquímico medido. \\
\texttt{7e7db277} & Emisiones de SO2 y formación de partículas. & \texttt{None} en versiones conservadoras. & PB9 si la evidencia particulada es explícita. & Añadir inferencias químicas permitidas cuando el abstract las formula. \\
\texttt{6adf36ff} & Metales pesados y contaminantes tóxicos. & \texttt{None} por vocabulario PB8 incompleto. & PB8 como entidad novedosa/contaminante. & Expandir PB8 con heavy metals, Pb, Cu, Mn y sinónimos. \\
\texttt{9153f4dc} & Simulación meteorológica sin cambio climático activo. & PB1 por vocabulario atmosférico. & \texttt{None} si solo es weather modelling. & No confundir meteorología operacional con PB1. \\
\texttt{5b69b4b1} & Base de datos sobre ozono. & \texttt{None} por filtro excesivo. & PB3 si la capa de ozono estratosférico es el objeto. & Separar ozono estratosférico de calidad del aire troposférica. \\
\bottomrule
\end{tabular}
\end{center}

\subsubsection{9.1 Falsos positivos por ground truth discutible}

Algunos casos donde el humano etiquetó \texttt{None} y el modelo predijo una PB no son necesariamente errores obvios del modelo. Por ejemplo, papers sobre PM2.5, black carbon, aerosol optical depth o aerosoles urbanos pueden haber sido marcados como \texttt{None} en la validación inicial, pero desde el marco PB9 sí contienen variables activas. Esto sugiere ruido de ground truth y falta de doble anotación independiente.

\subsubsection{9.2 Confusión macro-driver vs micro-impact}

Varios fallos ocurren cuando el abstract menciona cambio climático como driver, pero el fenómeno medido pertenece a otra PB. Ejemplos:

\begin{itemize}
\item Warming como causa, pero impacto real en genes microbianos o ecosistemas: PB7 puede ser primaria.
\item Clima como escenario, pero variable operacional de agua: PB5.
\item Dust o aerosoles como vector, pero medición principal de nutrientes: PB4.
\end{itemize}

V4 mejora este problema con el principio de \emph{operational object}.

\subsubsection{9.3 Anclaje en señales auxiliares}

La cascada introdujo una fuente nueva de error: si el scorer o el extractor proponen una PB incorrecta, el juez 14B puede anclarse en ella. Esto explica que la primera cascada rindiera peor que v4. La mitigación fue reordenar el prompt, mover señales auxiliares al final y declararlas explícitamente como falibles.

\subsubsection{9.4 Vocabulario incompleto en PB8 y PB1}

PB8 tiene baja cobertura cuando aparecen contaminantes como heavy metals no recogidos en la lista de novel entities. PB1 también puede perder paleoclimate proxies o indicadores indirectos si no aparecen términos clásicos de climate change. La mejora natural es expandir \texttt{pb\_reference.csv} mediante minería de errores.

\subsubsection{9.5 Clases con bajo soporte}

PB5 y PB6 tienen soporte muy bajo en la validación final (\texttt{n=4} cada una en varias evaluaciones). No se deben sacar conclusiones fuertes por PB individual con tan pocos ejemplos. Para estas clases conviene informar macro-F1, pero también advertir que la incertidumbre estadística es alta.

\subsection{10. Comparación con modelos embedding/BERT}

Además de LLMs, el proyecto evaluó baselines de embeddings y modelos tipo BERT/SciBERT/SPECTER en \texttt{nlp/bert\_finetuning/outputs/}. La mejor referencia no generativa actual es \texttt{baseline\_semantic\_tfidf} top-1, con:

\begin{itemize}
\item Micro-F1: \textbf{0.484}
\item Macro-F1: \textbf{0.439}
\item Jaccard: \textbf{0.367}
\item Exact Match: \textbf{0.280}
\item LRAP: \textbf{0.730}
\end{itemize}

\begin{center}
\scriptsize
\setlength{\tabcolsep}{4pt}
\renewcommand{\arraystretch}{1.15}
\begin{tabular}{p{0.25\linewidth}p{0.12\linewidth}p{0.12\linewidth}p{0.12\linewidth}p{0.12\linewidth}p{0.12\linewidth}}
\toprule
Modelo / estrategia & Micro-F1 & Macro-F1 & Jaccard & Exact Match & LRAP \\
\midrule
Baseline lexical threshold/delta & 0.485 & 0.402 & -- & 0.473 & -- \\
Baseline semantic TF-IDF top-1 & 0.484 & 0.439 & 0.367 & 0.280 & 0.730 \\
BERT top-1 & 0.349 & 0.341 & 0.260 & 0.193 & 0.638 \\
SciBERT top-2 & 0.348 & 0.373 & -- & 0.053 & 0.617 \\
SPECTER threshold/delta & 0.405 & 0.399 & -- & 0.313 & 0.648 \\
\bottomrule
\end{tabular}
\end{center}

La lectura técnica de esta tabla es que los modelos embedding funcionan bien como recuperadores de similitud semántica, pero no resuelven completamente la decisión normativa. La frontera entre \texttt{None} y PB activa depende de exclusiones, jerarquía causal y objeto medido, no solo de cercanía vectorial. Por eso el diseño final no reemplaza los embeddings por LLMs de forma dogmática: toma la parte fuerte de cada familia. El scorer conserva la transparencia léxica; el LLM aporta lectura contextual y razonamiento de abstención.

También se evaluaron BERT, RoBERTa, SciBERT y SPECTER. SPECTER con threshold/delta alcanzó Micro-F1 \textbf{0.405} y Macro-F1 \textbf{0.399}, y el baseline lexical threshold/delta alcanzó Micro-F1 \textbf{0.485}. Estos resultados muestran dos cosas:

\begin{enumerate}
\item Los métodos transparentes son sorprendentemente competitivos en vocabulario científico.
\item Los LLMs aportan una capa que los embeddings no capturan bien: capacidad de rechazar usos contextuales, explicar la decisión y distinguir objeto operacional frente a motivación.
\end{enumerate}

La solución final no abandona las señales léxicas; las incorpora como scorer determinista dentro de la cascada.

\subsection{11. Limitaciones}

Las limitaciones principales son:

\begin{itemize}
\item \textbf{Ground truth limitado y ruidoso}: 150 filas evaluables no bastan para estimar con precisión clases minoritarias. Hay duplicados y casos discutibles.
\item \textbf{Sin acuerdo inter-anotador}: falta medir cuán consistente sería la clasificación humana entre varios revisores.
\item \textbf{Dependencia del abstract}: algunos papers pueden requerir full text para decidir correctamente.
\item \textbf{Auto-confianza no calibrada}: \texttt{High/Medium/Low} es útil como señal cualitativa, pero no equivale a probabilidad real.
\item \textbf{Prompts largos}: mejoran contexto, pero aumentan coste y pueden introducir lost-in-the-middle.
\item \textbf{Riesgo de overfitting al subset}: cada iteración se evaluó sobre el mismo conjunto pequeño; la validación externa sería el siguiente paso.
\item \textbf{Salida multi-etiqueta infraexplotada}: la cascada predice secundarias, pero gran parte del análisis sigue usando la primera PB humana como referencia principal.
\end{itemize}

Estas limitaciones no invalidan el sistema; delimitan su modo de uso. El clasificador no debe presentarse como autoridad final automática, sino como herramienta de priorización y auditoría para curación humana.

\subsection{12. Valor del enfoque}

El valor del componente LLM no es únicamente subir una métrica. Aporta:

\begin{itemize}
\item \textbf{Escalabilidad}: permite pasar de revisión manual de decenas de papers a procesamiento de miles.
\item \textbf{Auditabilidad}: cada decisión incluye razonamiento y evidencias intermedias.
\item \textbf{Flexibilidad científica}: las reglas PB pueden editarse sin reentrenar un modelo.
\item \textbf{Privacidad}: la inferencia local evita enviar abstracts institucionales a APIs externas.
\item \textbf{Human-in-the-loop}: casos con baja confianza o discrepancias pueden priorizarse para revisión.
\item \textbf{Reutilización}: la arquitectura se puede adaptar a SDGs, áreas estratégicas UPV u otros marcos de impacto.
\end{itemize}

Para el producto UPV-EARTH, el uso realista es un sistema de screening: el modelo propone PBs, justifica la decisión y permite que un experto revise casos límite. Esto reduce carga manual y mejora trazabilidad institucional, pero evita presentar el LLM como juez final infalible.

\subsection{13. Encaje con la rúbrica de modelos}

La parte LLM encaja especialmente en los criterios de modelado, evaluación y mejora iterativa de la rúbrica. La siguiente tabla puede usarse como puente directo entre el trabajo técnico y lo que se espera justificar en el report final.

\begin{center}
\scriptsize
\setlength{\tabcolsep}{4pt}
\renewcommand{\arraystretch}{1.18}
\begin{tabular}{p{0.20\linewidth}p{0.31\linewidth}p{0.31\linewidth}p{0.12\linewidth}}
\toprule
Criterio defendible & Evidencia en el proyecto & Cómo se argumenta en la memoria & Artefactos \\
\midrule
Elección del modelo & Comparación entre Qwen 14B, Gemma 26B y baselines BERT/TF-IDF. & Qwen se elige por equilibrio entre calidad, latencia, coste local y estabilidad JSON. & Runners, tablas de métricas, logs. \\
Diseño metodológico & Prompts con reglas PB, activación/exclusión, objeto operacional y abstención. & El prompt es parte del modelo porque codifica la definición científica de etiqueta. & Sección 5.6, \texttt{pb\_reference.csv}. \\
Evaluación cuantitativa & Top-1, PB-only, None-only, Top-2 global, Hit@K PB-only, Exact Match, Jaccard, Hamming y positivity bias. & No se optimiza una métrica única; se evalúa precisión, cobertura y rigor. & CSVs de inferencia y notebooks de análisis. \\
Iteración justificada & v1 empeora por ejemplos rígidos; v2 aumenta recall pero sesga; v3 recupera rigor; v4 mejora global. & Cada cambio responde a un fallo observado, no a tuning arbitrario. & Tabla de iteraciones y análisis pareado. \\
Arquitectura avanzada & Cascada con extractor, scorer, router, juez y critic. & El sistema final prioriza auditabilidad y revisión humana además de accuracy. & \texttt{pipeline\_agentes.py}, figuras 1--2. \\
Análisis crítico & Taxonomía de errores, casos cualitativos, límites de ground truth. & Se reconocen ruido, bajo soporte de clases y necesidad de reanotación. & Figuras 4, 8 y tabla de casos. \\
Impacto práctico & Screening escalable sobre corpus institucional ampliado. & El modelo reduce carga manual y ofrece trazabilidad sin sustituir juicio experto. & Outputs y propuesta human-in-the-loop. \\
\bottomrule
\end{tabular}
\end{center}

\subsection{14. Próximos pasos técnicos}

Los siguientes pasos más defendibles son:

\begin{enumerate}
\item \textbf{Reanotación doble del ground truth}: al menos dos revisores independientes y cálculo de acuerdo.
\item \textbf{Evaluación multi-etiqueta completa}: usar todas las columnas \texttt{1stpb}, \texttt{2ndpb}, \texttt{3rdpb}, no solo primaria.
\item \textbf{Expansión de \texttt{pb\_reference.csv} desde errores}: heavy metals para PB8, paleoclimate proxies para PB1, términos de conservación para PB7.
\item \textbf{Critic ampliado}: permitir revisión no solo de \texttt{None -> PBx}, sino también de casos de baja confianza o desacuerdo fuerte con \texttt{kw\_top}.
\item \textbf{Calibración de abstención}: introducir estado \texttt{Uncertain} para enviar a humano.
\item \textbf{Destilación de razonamientos Gemma $\rightarrow$ Qwen}: usar Gemma como auditor offline para mejorar prompts o crear datos débiles.
\item \textbf{Evaluación externa}: aplicar el sistema a un subconjunto nuevo del corpus institucional ampliado.
\end{enumerate}

\subsection{15. Archivos de trazabilidad}

Los artefactos principales están versionados en el repositorio:

\begin{itemize}
\item Reglas PB: \texttt{corpus\_PB/data/pb\_reference.csv}
\item Ground truth: \texttt{nlp/llm/outputs/ground\_truth/validacion\_real.csv}
\item Zero-shot runner: \texttt{nlp/llm/runners/qwen\_zeroshot.py}
\item V2 runner: \texttt{nlp/llm/runners/qwen\_fewshot\_v2.py}
\item V3 runner/notebook: \texttt{nlp/llm/runners/qwen\_fewshot\_v3.ipynb}
\item V4 runner/notebook: \texttt{nlp/llm/runners/qwen\_fewshot\_v4.ipynb}
\item Cascada: \texttt{nlp/llm/runners/pipeline\_agentes.py}
\item Critic post-hoc: \texttt{nlp/llm/runners/apply\_critic.py}
\item Análisis v1-v4: \texttt{nlp/llm/analysis/comparacion\_fewshot\_v1\_v2\_v3\_v4\_principle.ipynb}
\item Análisis cascada: \texttt{nlp/llm/analysis/pipeline\_cascada\_analysis.ipynb}
\item Reporte cascada: \texttt{docs/report/multi\_agent\_pipeline\_report.md}
\item Figuras listas: \texttt{docs/report/figures/}
\end{itemize}

\subsection{16. Texto breve para cerrar la sección en el paper}

En síntesis, la evolución del componente LLM muestra una mejora incremental real y bien diagnosticada. El sistema pasó de prompts ingenuos con sesgo de positividad a un clasificador principle-driven basado en reglas explícitas de activación/exclusión, salida JSON y auditoría por razonamiento. Qwen 2.5 14B v4 alcanzó el mejor Top-1 single-model (72.1\%), mientras que la cascada multi-agente final alcanzó 71.3\% Top-1, 73.3\% Top-2 global y 74.7\% Hit@K PB-only, añadiendo trazabilidad mediante extracción estructurada, scoring determinista y critic asimétrico. El resultado final no debe interpretarse como una sustitución completa del juicio experto, sino como una herramienta escalable y auditable para priorizar, explicar y revisar la contribución científica de la UPV al marco de Planetary Boundaries.
